\uuid{IchG}
\exo7id{4529}
\titre{exo7 4529}
\auteur{quercia}
\organisation{exo7}
\datecreate{2010-03-14}
\isIndication{false}
\isCorrection{true}
\chapitre{Suite et série de fonctions}
\sousChapitre{Autre}
\cpge{mpsi-mp2i.nombres-reels-suites.d, pcsi.nombres-reels-suites.c, ptsi.nombres-reels-suites.c, mpsi-mp2i.limites-continuite.a, pcsi.limites-continuite.a, ptsi.limites-continuite.a, bcpst1.limites-continuite.a, bcpst1.fonctions-reelles-usuelles.a, mp-mpi.series-numeriques-vectorielles.a, pt.series-numeriques.b, mp-mpi.series-entieres.d, pc.series-entieres.c, psi.series-entieres.c, pt.series-entieres.c, bcpst2.series-reelles, bcpst2.integrales-generalisees}
\module{Analyse}
\niveau{L2}
\difficulte{}

\contenu{
\texte{
Soit~$f(a) = \sum_{n=0}^\infty e^{-a^2n^2}$ sous réserve de convergence $(a\in\R$).
}
\begin{enumerate}
    \item \question{Domaine de définition de~$f$~?}
\reponse{$\R^*$.}
    \item \question{Limite de~$af(a)$ quand $a\to0$~?}
\reponse{CSI~: $\frac{\sqrt\pi}{2a} =  \int_{x=0}^\infty e^{-a^2x^2}\,d x \le f(a) \le  \int_{x=0}^\infty e^{-a^2x^2}\,d x + 1 = \frac{\sqrt\pi}{2a} + 1$. Donc $af(a)\to \frac{\sqrt\pi}2$ lorsque $a\to0^+$.}
    \item \question{Limite de~$f(a)$ quand $a\to+\infty$~?}
\reponse{TCM~: $f(a)\to 1$ lorsque $a\to+\infty$.}
\end{enumerate}
}
